A mathematical model of language contact from previous literature \parencite{Kauhanen2022} predicts a critical threshold for contact-induced change: the proportion of second-language (L2) learners present in a population of speakers must exceed this threshold in order for simplification to occur. In this chapter, I investigate to what extent the model's predictions remain valid when key idealizing assumptions are dropped. Rather than working in the model's deterministic limit in which analytical results are available, the model is here studied using agent-based simulations, using realistically sized speech communities and locally structured interaction patterns. It is found that the outcomes of the agent-based simulations closely match those predicted by the deterministic limit, up to minute deviances due to finite-size effects. Additionally, the interactions of several key model parameters are investigated, and the two specific case studies of \textcite{Kauhanen2022} are revisited using the agent-based methodology.
